% ch9_a 第9章 §9.1 (书页 354–366 / p369–p381)
% 第9章 MEAN-FIELD THEORY OF SPIN LIQUIDS AND QUANTUM ORDER

\chapter{自旋液体与量子序的平均场理论}

\begin{keypoints}
\item 用从属玻色子方法（或投影构造）建立自旋液体态的平均场理论。
\item 规范相互作用的重要性，以及如何应用（含规范相互作用的）平均场理论来研究自旋液体的物理性质。
\item 自旋液体的表征，以及拓扑序与量子序的概念。
\end{keypoints}

在本章中，我们将为自旋液体态建立一个平均场理论。这里所谓"自旋液体态"，是指具有自旋旋转对称性、且每个元胞含\emph{奇数}个电子的绝缘体。通常，每个元胞含奇数个电子的态具有半填充能带，是导体。因此，自旋液体——如果它们存在的话——是非常奇特的态。自旋液体如此奇异，以致许多人根本不相信它们能够存在。事实上，即使到了现在，我们也还不知道有哪个自旋哈密顿量能被证明可以可靠地给出自旋液体基态。

对于那些相信自旋液体存在的人，他们必须接受或相信自旋液体的下列奇特（或者说迷人）的性质。（i）自旋液体中的激发总是携带分数量子数，例如电中性的自旋 $1/2$，这从电子的任何集合出发都不可能得到；有时激发甚至可以携带分数统计。（ii）不同的自旋液体无法靠它们的对称性质来区分；自旋液体是量子有序态的例子。（iii）有能隙自旋液体的基态总是具有拓扑简并，这种简并不能与任何对称性相联系。（iv）自旋液体总是包含某种规范涨落。

在本章中，我们将假定自旋液体确实存在，并为自旋液体建立平均场理论。平均场理论使我们能够理解自旋液体的某些物理性质。在计入重要的平均场涨落之后，我们将证明，某些平均场态对这些涨落是稳定的，因而代表真实的自旋液体态。由于自旋液体是具有非平凡拓扑序/量子序的典型态，我们将以自旋液体为载体来建立量子序的理论。

\section{量子自旋液体态的投影构造}

在这一节中，我们将使用从属玻色子方法（slave-boson approach，或称投影构造（projective construction））（Baskaran et al., 1987; Affleck et al., 1988; Baskaran and Anderson, 1988; Dagotto et al., 1988; Wen and Lee, 1996; Senthil and Fisher, 2000）来构造二维自旋液体。Baskaran 和 Anderson（1988）在从属玻色子方法中发现的规范结构，对我们理解强关联自旋液体起着关键作用。除从属玻色子方法之外，还可以用另一类投影构造——从属费米子方法或 Schwinger 玻色子方法——来构造自旋液体态（Arovas and Auerbach, 1988; Read and Sachdev, 1991; Sachdev and Park, 2002）。由于从属费米子方法只能给出具有有限能隙的自旋液体，这里我们将集中于从属玻色子方法。

\subsection{自旋液体态的平均场理论}

\begin{keypoints}
\item 自旋液体的平均场理论可以由投影构造得到。
\item 平均场理论以及 $\pi$ 通量相的平均场拟设（ansatz）。
\item $a_0(\pmb{i})$ 的规范涨落实施约束条件。
\end{keypoints}

让我们考虑二维正方格子上的 Hubbard 模型（式 (5.5.1)）。在半填充时，Hubbard 模型化为如下的海森堡模型：
\begin{equation*}
H=\sum_{\langle ij\rangle}J_{ij}\,\pmb{S}_i\cdot\pmb{S}_j
\tag{9.1.1}
\end{equation*}

我们想指出，上述自旋 $1/2$ 模型也可以看作一个硬核玻色子（hard-core boson）模型，其中 $|\downarrow\rangle$ 对应空格点，$|\uparrow\rangle$ 对应被一个玻色子占据的格点。上述哈密顿量很难求解，于是我们用平均场近似来理解它的物理性质。在平均场近似中，我们把 $\pmb{S}_i$ 之一换成其量子平均值 $\langle\pmb{S}_i\rangle$，得到如下平均场哈密顿量：
\begin{equation*}
H_{\text{mean}}=\sum_{\langle ij\rangle}J_{ij}\left(\langle\pmb{S}_i\rangle\cdot\pmb{S}_j+\pmb{S}_i\cdot\langle\pmb{S}_j\rangle-\langle\pmb{S}_i\rangle\cdot\langle\pmb{S}_j\rangle\right)
\end{equation*}

平均场哈密顿量 $H_{\text{mean}}$ 容易求解，我们可以得到 $H_{\text{mean}}$ 的基态，即 $|\Phi_{\text{mean}}\rangle$。我们唯一需要做的，是仔细选取 $\langle\pmb{S}_i\rangle$ 的值，使它们满足所谓的自洽方程
\begin{equation*}
\langle\pmb{S}_i\rangle=\langle\Phi_{\text{mean}}|\pmb{S}_i|\Phi_{\text{mean}}\rangle
\end{equation*}

这就是自旋系统的标准平均场方法。

上述标准平均场方法的问题在于，它只能用于研究有序的自旋态，因为我们从一开始就假定了 $\langle\pmb{S}_i\rangle\neq0$。长期以来，自旋液体的平均场理论似乎是不可能建立的。1987 年，我们终于找到了一个奇特的技巧——从属玻色子方法\footnote{从属玻色子（slave boson）和从属费米子（slave fermion）是两个非常奇怪且令人困惑的术语。从属玻色子方法里没有玻色子，从属费米子方法里没有费米子。这就是我更愿意用"投影构造"来称呼这两种方法的原因。}——来做这件事（Baskaran et al., 1987）。为了得到自旋液体的平均场基态，我们引入自旋子（spinon）算符 $f_{i\alpha}$，$\alpha=1,2$，它们是自旋 $1/2$ 的电中性算符。自旋算符 $\pmb{S}_i$ 表示为
\begin{equation*}
\pmb{S}_i=\frac{1}{2}f^{\dagger}_{i\alpha}\sigma_{\alpha\beta}f_{i\beta}
\tag{9.1.2}
\end{equation*}

用自旋子算符表示，哈密顿量 (9.1.1) 可以改写为
\begin{equation*}
H=\sum_{\langle ij\rangle}-\frac{1}{2}J_{ij}f^{\dagger}_{i\alpha}f_{j\alpha}f^{\dagger}_{j\beta}f_{i\beta}+\sum_{\langle ij\rangle}J_{ij}\left(\frac{1}{2}n_i-\frac{1}{4}n_in_j\right)
\tag{9.1.3}
\end{equation*}

这里我们用了 $\sigma_{\alpha\beta}\cdot\sigma_{\alpha'\beta'}=2\delta_{\alpha\beta'}\delta_{\alpha'\beta}-\delta_{\alpha\beta}\delta_{\alpha'\beta'}$，而 $n_i$ 是格点 $i$ 上的费米子数。式 (9.1.3) 中的第二项是常数，在下面的讨论中将略去。注意，式 (9.1.3) 的希尔伯特空间每个格点有四个态，比每格点只有两个态的式 (9.1.1) 的希尔伯特空间要大。式 (9.1.1) 与式 (9.1.3) 之间的等价只在每个格点恰好有一个费米子的子空间中才成立。因此，要用式 (9.1.3) 来描写自旋态，我们需要施加约束（Baskaran et al., 1987; Baskaran and Anderson, 1988）
\begin{equation*}
f^{\dagger}_{i\alpha}f_{i\alpha}=1,\qquad f_{i\alpha}f_{i\beta}\epsilon_{\alpha\beta}=0
\tag{9.1.4}
\end{equation*}

第二个约束实际上是第一个约束的推论。

零阶平均场基态可以通过如下近似得到。首先，我们用约束 (9.1.4) 的基态平均值
\begin{equation*}
\langle f^{\dagger}_{i\alpha}f_{i\alpha}\rangle=1
\tag{9.1.5}
\end{equation*}

来代替约束本身。这一约束可以通过在哈密顿量中引入一个依赖格点且不依赖时间的拉格朗日乘子 $a_0(\pmb{i})(f^{\dagger}_{i\alpha}f_{i\alpha}-1)$ 来实施。其次，我们把算符 $f^{\dagger}_{i\alpha}f_{j\alpha}$ 换成其基态期望值 $\chi_{ij}$，同样忽略它们的涨落。这样，我们便得到零阶平均场哈密顿量
\begin{equation*}
H_{\text{mean}}=\sum_{\langle ij\rangle}-\frac{1}{2}J_{ij}\left[\left(f^{\dagger}_{i\alpha}f_{j\alpha}\chi_{ji}+\text{h.c}\right)-|\chi_{ij}|^{2}\right]+\sum_{i}a_0(\pmb{i})\left(f^{\dagger}_{i\alpha}f_{i\alpha}-1\right)
\tag{9.1.6}
\end{equation*}

式 (9.1.6) 中的 $\chi_{ij}$ 必须满足自洽条件
\begin{equation*}
\chi_{ij}=\langle f^{\dagger}_{i\alpha}f_{j\alpha}\rangle
\tag{9.1.7}
\end{equation*}

而依赖格点的化学势 $a_0(\pmb{i})$ 则选得使平均场基态满足式 (9.1.5)。为方便起见，我们将把 $\chi_{ij}$ 称为自旋液体态的"拟设"。

设 $\bar{\chi}_{ij}$ 是式 (9.1.7) 的一个解，$\bar{a}(\pmb{i})$ 是式 (9.1.5) 的一个解。这样的拟设对应于平均场基态。由于 $\chi_{ij}$ 和 $a_0$ 在自旋旋转变换下不变，平均场基态在自旋旋转变换下不变。如果我们忽略 $\bar{\chi}_{ij}$ 和 $\bar{a}(\pmb{i})$ 的涨落，则平均场态周围的激发由下式描写：
\begin{equation*}
H_{\text{mean}}=\sum_{\langle ij\rangle}-\frac{J_{ij}}{2}\left[\left(f^{\dagger}_{i\alpha}f_{j\alpha}\bar{\chi}_{ji}+\text{h.c}\right)-|\bar{\chi}_{ij}|^{2}\right]+\sum_{i}\bar{a}_0(\pmb{i})\left(f^{\dagger}_{i\alpha}f_{i\alpha}-1\right)
\end{equation*}

我们看到，零阶平均场理论中的激发是由 $f_{i\alpha}$ 描写的自由自旋子。自旋子是自旋 $1/2$ 的中性费米子。令人惊叹的是，费米子激发竟然能从一个纯玻色的模型——式 (9.1.1)——中演生出来。

现在的问题是：我们是否应该相信这个平均场结果？我们是否应该相信具有自旋 $1/2$ 中性费米子激发的自旋液体能够存在？检验的办法之一是把平均场拟设周围的涨落包括进来，看看涨落是否会改变平均场结果。因此，下面我们将考虑涨落的效应。

首先，我们想指出，如果当初把 $a_0$ 的涨落（即时间依赖性）包括进来，约束 (9.1.5) 就会变回原来的约束 (9.1.4)。为看清这一点，让我们考虑 $H_{\text{mean}}$ 的路径积分表述：
\begin{gather*}
Z=\int\mathcal{D}f\,\mathcal{D}[a_0(\pmb{i})]\,\mathcal{D}\chi_{ij}\,\mathrm{e}^{\mathrm{i}\int\mathrm{d}t\,(\mathcal{L}-\sum_{i}a_0(\pmb{i},t)(f^{\dagger}_if_i-1))}\\
\mathcal{L}=\sum_{i}f^{\dagger}_i\mathrm{i}\partial_tf_i-\sum_{\langle ij\rangle}-\frac{1}{2}J_{ij}\left[\left(f^{\dagger}_{i\alpha}f_{j\alpha}\chi_{ji}+\text{h.c}\right)-|\chi_{ij}|^{2}\right]
\tag{9.1.8}
\end{gather*}

我们看到，对依赖时间的 $a_0(\pmb{i},t)$ 作积分将产生一个约束
\begin{equation*}
\prod_{i,t}\delta\left(f^{\dagger}_i(t)f_i(t)-1\right)
\end{equation*}

它在每个格点 $i$、每个时刻 $t$ 都得到实施。我们还注意到，对 $\chi_{ij}(t)$ 作积分将复原原来的哈密顿量 (9.1.3)。因此，式 (9.1.8) 是自旋模型 (9.1.3) 的一个严格表示。$\chi_{ij}$ 和 $a_0(\pmb{i})$ 中的涨落描写平均场基态之上的集体激发。

这里 $\chi_{ij}$ 有两种涨落：振幅涨落和相位涨落。振幅涨落具有有限能隙，在我们的讨论中并不关键。所以这里我们只考虑平均场拟设 $\bar{\chi}_{ij}$ 周围的相位涨落 $a_{ij}$：
\begin{equation*}
\chi_{ij}=\bar{\chi}_{ij}\mathrm{e}^{-\mathrm{i}a_{ij}}
\tag{9.1.9}
\end{equation*}

把上述相位涨落和 $a_0$ 的涨落包括进来后，平均场哈密顿量变为
\begin{equation*}
H=\sum_{\langle ij\rangle}-J_{ij}\left(f^{\dagger}_{i\alpha}f_{j\alpha}\bar{\chi}_{ji}\mathrm{e}^{-\mathrm{i}a_{ji}}+\text{h.c}\right)-\sum_{i}a_0(\pmb{i})\left(f^{\dagger}_{i\alpha}f_{i\alpha}-1\right)
\tag{9.1.10}
\end{equation*}

我们想把式 (9.1.10) 称为一阶平均场哈密顿量。由式 (9.1.10) 可见，$a_0$ 和 $a_{ij}$ 所描写的涨落正是一个 $U(1)$ 格点规范场（见式 (6.4.9)）（Baskaran and Anderson, 1988; Lee and Nagaosa, 1992）。该哈密顿量在规范变换
\begin{equation*}
a_{ij}\rightarrow a_{ij}+\theta_i-\theta_j,\qquad f_i\rightarrow f_i\,\mathrm{e}^{\mathrm{i}\theta_i}.
\tag{9.1.11}
\end{equation*}

之下不变。

因此，在一阶平均场理论中，激发由耦合到 $U(1)$ 规范场的自旋子描写（而非零阶平均场理论中的自由自旋子）。

在通常的平均场理论中，我们对相互作用的哈密顿量作近似以简化问题，希尔伯特空间不因平均场近似而改变。投影构造（或从属玻色子方法）在这一点上非常不同：不仅哈密顿量改变了，希尔伯特空间也改变了。这使得由投影构造得到的平均场结果不仅定量上不正确，而且定性上也不正确。例如，平均场基态 $|\Psi_{\text{mean}}\rangle$（式 (9.1.6) 中 $H_{\text{mean}}$ 的基态）甚至不是一个合法的自旋波函数，因为有些格点上可能没有费米子或有两个费米子。然而，不要放弃。在通常的平均场理论中，我们可以通过把平均场态周围的涨落包括进来，在定量上改进近似。在上面我们看到，在投影构造中，我们可以通过把平均场态周围的规范涨落（$a_0,a_{ij}$）包括进来，在定性上改进近似并复原原来的希尔伯特空间。因此，要想从投影构造得到哪怕是定性的结果，至少把平均场态周围的规范涨落包括进来是至关重要的。换句话说，我们在自旋液体的投影构造中把一个自旋切成了两半；把这两半重新粘合起来，以得到正确的物理结果，是十分重要的。

按照零阶平均场理论，自旋液体中的低能激发是自由自旋子。由上面的讨论我们看到，这样的结果是不正确的。按照一阶平均场理论，自旋子通过规范涨落相互作用。规范相互作用可能会剧烈地改变自旋子的性质。我们稍后将回到这个问题。

\subsection*{问题 9.1.1}

证明式 (9.1.3)。

\subsection{信还是不信}

\begin{keypoints}
\item 一阶平均场理论的解禁闭相导致新一类物质态——量子有序态。
\item 向物理自旋波函数的投影，以及 $U(1)$ 规范结构的含义。
\item 作为纠缠之涨落的演生规范玻色子与费米子。
\end{keypoints}

按照一阶平均场理论，平均场基态周围的涨落由规范场和费米子场描写。别忘了，我们原来的模型只是一个相互作用的纯玻色自旋模型。一个纯玻色的模型怎么会有一个由规范场和费米子场描写的有效理论呢？这简直难以置信。让我们检查一下自己是怎么走到这一步的。我们先把玻色的自旋算符拆成两个费米子算符（自旋子算符）的乘积，然后引入一个规范场把自旋子重新粘合成玻色的自旋。从这个角度看，一阶平均场理论好像只是一个假的理论：规范玻色子和费米子似乎都是假的，到头来我们拥有的只不过是玻色自旋涨落。

然而，我们不应过快地抛弃一阶平均场理论。如果规范场处于禁闭相（见 6.4.3 节），它其实能够复现上述玻色自旋涨落的图像。在禁闭相中，自旋子之间通过线性势相互作用，永远不能在低能下作为准粒子出现；规范玻色子在禁闭相中有很大的能隙，不出现在低能谱中。唯一的低能激发是自旋子对，它们对应于玻色自旋涨落。所以一阶平均场理论也许没有用，但它并没有错。它能够给出与常识一致的图像（尽管绕了很远的弯路）。

另一方面，当规范场处于解禁闭相时，一阶平均场理论也能给出有悖常识的图像。在这种情况下，自旋子和规范玻色子将作为良定义的准粒子出现。问题是：我们相信解禁闭相的图像吗？我们相信从纯玻色模型中能够演生出规范玻色子和费米子吗？显然，上面概述的投影构造太过形式化，很难让大多数人相信如此激进的结果。

我不得不说，把自旋切成两半再粘回去这套做法让很多人望而却步。很难看出这种形式操作能带来什么新的物理洞见和结果。如果这条路径真的给出了什么新结果，许多人也会把它们归咎于这一奇怪构造的人为假象，而不是自旋液体的真实物理性质。的确，投影构造非常形式化，但它也是一件超越时代的礼物。我们现在知道，投影构造真正所做的，是产生一个弦网凝聚态。弦网凝聚（string-net condensation）给出一种新型的关联态，它具有非平凡的量子序（见第 10 章）。所以，请相信我：只要恰当地计入规范涨落的效应，投影构造给出的惊人结果是可以信赖的。

这里让我们暂且不用弦网凝聚的图像，试着理解投影构造背后的物理图像。首先，我们需要理解平均场拟设 $\chi_{ij}$ 是如何与每个格点恰好有一个费米子的物理自旋波函数相联系的。我们知道，平均场波态 $|\Psi^{(\chi_{ij})}_{\text{mean}}\rangle$（$H_{\text{mean}}$ 的基态）不是自旋系统的合法波函数，因为它未必每格点有一个费米子。为了把它联系到物理自旋波函数，我们需要把 $a_0$ 的涨落包括进来，以实施每格点一个费米子的约束。有了这一理解，我们就可以把平均场态投影到每格点一个费米子的子空间上，得到自旋系统的一个合法波函数 $\Psi_{\text{spin}}(\{\alpha_i\})$：
\begin{equation*}
\Psi^{(\chi_{ij})}_{\text{spin}}(\{\alpha_i\})=\langle 0_f|\prod_{i}f_{i\alpha_i}|\Psi^{(\chi_{ij})}_{\text{mean}}\rangle.
\tag{9.1.12}
\end{equation*}

其中 $|0_f\rangle$ 是没有 f 费米子的态，即 $f_{i\alpha}|0_f\rangle=0$。式 (9.1.12) 把平均场拟设（连同它的涨落）与物理自旋波函数联系起来，使我们能够理解平均场拟设和平均场涨落的物理意义。

例如，投影 (9.1.12) 赋予规范变换 (9.1.11) 以物理意义。由规范变换
\begin{equation*}
\tilde{\chi}_{ij}=\mathrm{e}^{\mathrm{i}\theta_i}\chi_{ij}\mathrm{e}^{-\mathrm{i}\theta_j}
\tag{9.1.13}
\end{equation*}

联系的两个平均场拟设 $\chi_{ij}$ 与 $\tilde{\chi}_{ij}$ 给出同一个投影后的自旋态
\begin{equation*}
\Psi^{(\tilde{\chi}_{ij})}_{\text{spin}}(\{\alpha_i\})=\mathrm{e}^{\mathrm{i}\sum_i\theta_i}\Psi^{(\chi_{ij})}_{\text{spin}}(\{\alpha_i\})
\tag{9.1.14}
\end{equation*}

对于 $\chi_{ij}$ 的不同取值，式 (9.1.6) 中 $H_{\text{mean}}$ 的基态对应于不同的平均场波函数 $|\Psi^{(\chi_{ij})}_{\text{mean}}\rangle$；投影之后，它们给出不同的物理自旋波函数 $\Psi^{(\tilde{\chi}_{ij})}_{\text{spin}}(\{\alpha_i\})$。因此，我们可以把 $\chi_{ij}$ 看作标记不同物理自旋态的标签。式 (9.1.14) 告诉我们，这个标签不是一对一的，而是多对一的。这一性质对我们理解 $\chi_{ij}$ 涨落的反常动力学性质十分重要。用多个标签标记同一个物理态，也使我们的理论按照 6.1.1 节中的定义成为一门规范理论。

让我们考虑多对一性质，或者说 $\chi_{ij}$ 的规范结构，如何影响它的动力学性质。如果 $\chi_{ij}$ 是物理态的一对一标签，那它就会像玻色超流体中的凝聚玻色子振幅 $\langle\phi(\pmb{x},t)\rangle$，或 SDW 态中的凝聚自旋矩 $\langle\pmb{S}_i(t)\rangle$ 一样，而 $\chi_{ij}$ 的涨落将对应于一种类似自旋波模式的玻色型模式。\footnote{更确切地说，自旋波模式对应于标量玻色子。局域序参量的涨落总是给出标量玻色子。}然而，$\chi_{ij}$ 的行为并不像局域序参量 $\langle\phi(\pmb{x},t)\rangle$ 和 $\langle\pmb{S}_i(t)\rangle$ 那样无冗余地标记物理态。作为多对一的标签，$\chi_{ij}$ 的某些涨落不改变物理态，因而是非物理的；这些涨落被称为纯规范涨落（pure gauge fluctuations）。对一个一般的涨落 $\delta\chi_{ij}$，其中一部分是物理的，另一部分是非物理的。$\chi_{ij}$ 的有效理论必须是规范不变的。这一点剧烈地改变了涨落的动力学性质。正是这一性质使 $\chi_{ij}$ 的涨落表现得像规范玻色子，与声模式和自旋波模式大不相同。

我们已经论证过，$\chi_{ij}$ 的相位涨落对应于规范涨落。投影构造 (9.1.12) 使我们能够得到与规范涨落 $a_{ij}$ 相应的物理自旋波函数：
\begin{equation*}
\Psi^{(a_{ij})}_{\text{spin}}=\langle 0|\prod_{i}f_{i\alpha_i}|\Psi^{(\bar{\chi}_{ij}\mathrm{e}^{\mathrm{i}a_{ij}})}_{\text{mean}}\rangle.
\end{equation*}

类似地，投影构造也使我们能够得到与一对自旋子激发相应的物理自旋波函数。我们从带有一对粒子–空穴激发的平均场基态出发，经过投影 (9.1.12)，便得到含有一对自旋子的物理自旋波函数：
\begin{equation*}
\Psi^{\text{spinon}}_{\text{spin}}(\pmb{i}_1,\alpha_1;\pmb{i}_2,\alpha_2)=\langle 0|\Big(\prod_{i}f_{i\alpha_i}\Big)f^{\dagger}_{\pmb{i}_1\lambda_1}f_{\pmb{i}_2\lambda_2}|\Psi^{(\bar{\chi}_{ij})}_{\text{mean}}\rangle.
\end{equation*}

如果你对这里给出的物理图像不满意，仍然不相信投影构造能够产生演生的规范玻色子和费米子，那么你可以直接跳到第 10 章。在第 10 章中，我们构造了几个能够用投影构造精确（或准精确）求解的自旋模型。这些模型具有演生的 $Z_2/U(1)$ 规范结构和费米子。这些精确可解模型揭示了演生规范玻色子和费米子的弦网起源。在本章其余部分，我们假定投影构造是有意义的，并研究它的推论。

%FIG{9.1}: {二聚体态由自旋单态对构成。把一个二聚体变成三重态便产生一个自旋 1 激发。两个分离开的自旋 $1/2$ 激发由一条由错位二聚体构成的弦相连。}

\subsection{二聚体态}

\begin{keypoints}
\item 二聚体态的一阶平均场理论具有禁闭型的 $U(1)$ 规范相互作用。物理激发中只出现自旋子对的束缚态。
\end{keypoints}

让我们先讨论一个具体的平均场基态——具有最近邻耦合的海森堡模型中的二聚体态（dimer state）（Majumdar and Ghosh, 1969; Affleck et al., 1987; Read and Sachdev, 1989）：
\begin{equation*}
H=J_1\sum_{i}\left(\pmb{S}_i\pmb{S}_{\pmb{i}+\pmb{x}}+\pmb{S}_i\pmb{S}_{\pmb{i}+\pmb{y}}\right).
\tag{9.1.15}
\end{equation*}

二聚体态由如下拟设描写：
\begin{equation*}
\begin{array}{l}
\chi_{i,i+\pmb{x}}=\dfrac{1}{2}\left(1+(-)^{i_x}\right)\\[6pt]
\text{其余}\;=0
\end{array}
\tag{9.1.16}
\end{equation*}

以及 $a_0(\pmb{i})=0$。可以验证，式 (9.1.6) 中 $H_{\text{mean}}$ 的平均场基态满足自洽方程 (9.1.7) 和约束 (9.1.5)。在平均场二聚体态中，自旋子能谱没有色散：
\begin{equation*}
E_k=\pm\frac{1}{2}J_1
\tag{9.1.17}
\end{equation*}

具有 $E_k=-\frac{1}{2}J_1$ 的价带被自旋子完全填满。自旋激发在二聚体态中有有限能隙。同样清楚的是，二聚体态破坏 $x$ 方向的平移对称性。

在零阶平均场理论中，二聚体基态之上的激发是自旋 $1/2$ 的自旋子。这一结果显然不对，因为从物理上看，二聚体态是由填满格子的自旋单态二聚体构成的（见图 9.1），基本激发对应于自旋三重态的二聚体，即玻色型的自旋 1 激发。在一阶平均场理论中，自旋子耦合到 $U(1)$ 规范场。在 $2+1$ 维，$U(1)$ 规范理论是禁闭的，因此单个自旋子不可观测，物理谱中只能出现玻色型的束缚态（携带整数自旋）。如果我们真的制造出两个分离开的自旋 $1/2$ 激发，那么由图 9.1 可见，它们将由一条错位二聚体构成的弦连接起来。这根弦在两个自旋 $1/2$ 激发之间导致线性禁闭相互作用，与一阶平均场理论给出的图像一致。我们看到，一阶平均场理论给出了定性上正确的结果。

%FIG{9.2}: {(a) $\pi$ 通量态的平均场拟设。沿箭头方向 $\chi_{ij}=\mathrm{i}\chi_1$。(b) $\pi$ 通量态中费米子的色散。价带是填满的。低能激发存在于价带与导带相切的四个费米点附近。}

\subsection{$\pi$ 通量态}

\begin{keypoints}
\item 投影之后，平均场 $\pi$ 通量相给出平移、旋转和宇称对称的自旋液体波函数。
\item $\pi$ 通量相的低能有效理论（一阶平均场理论）包含无能隙狄拉克费米子耦合到 $U(1)$ 规范场。这些狄拉克费米子携带自旋 $1/2$ 且电中性。
\item 稳定、边缘和不稳定平均场态的概念。
\item 一阶平均场理论只对稳定平均场态或涨落微弱的边缘平均场态可靠。$\pi$ 通量相不是这些平均场态之一。
\end{keypoints}

我们要研究的第二个平均场态，是最近邻海森堡模型的 $\pi$ 通量态（$\pi$-flux state）（Affleck and Marston, 1988; Kotliar, 1988）。$\pi$ 通量态由如下拟设给出（见图 9.2）：
\begin{equation*}
\chi_{i,i+\pmb{x}}=\mathrm{i}\chi_1,\qquad \chi_{i,i+\pmb{y}}=\mathrm{i}\chi_1(-)^{i_x},\qquad a_0(\pmb{i})=0.
\tag{9.1.18}
\end{equation*}

平均场哈密顿量 (9.1.6) 等价于每个方格带 $\pi$ 通量的电子跳跃问题，因为绕一个方格一周有 $\prod\chi_{ij}\propto\mathrm{e}^{\mathrm{i}\pi}$。$\pi$ 通量相中自旋子能谱为（见图 9.2）
\begin{equation*}
E_k=\pm J_1\chi_1\sqrt{\sin(k_x)^{2}+\sin(k_y)^{2}}
\tag{9.1.19}
\end{equation*}

其中 $-\frac{\pi}{2}<k_x<\frac{\pi}{2}$，$-\pi<k_y<\pi$。$\chi_1$ 的值可以由对平均场能量
\begin{equation*}
-\sum_{\pmb{k}}^{\prime}J_1\chi_1\sqrt{\sin(k_x)^{2}+\sin(k_y)^{2}}+J_1|\chi_1|^{2}N_{\text{site}}
\end{equation*}

求极小得到。我们发现
\begin{equation*}
\chi_1=\frac{1}{4}\int\frac{\mathrm{d}^{2}\pmb{k}}{(2\pi)^{2}}\sqrt{\sin(k_x)^{2}+\sin(k_y)^{2}}
\end{equation*}

除 $\pi$ 通量相之外，还有许多其他不同的平均场拟设局部地使平均场能量取极小。不过，对于只有最近邻耦合 $J_1$ 的海森堡模型，$\pi$ 通量相在平移不变的自旋液体中平均场能量最低。

让我们讨论 $\pi$ 通量相的一些物理性质。首先，考虑平均场态的对称性。我们知道，平均场理论旨在描写一个自旋态，其波函数为 $\Psi_{\text{spin}}(\{\alpha_i\})$，其中 $\alpha_i=\pm 1/2$ 是格点 $i$ 上 $S^{z}$ 的取值。当我们说"平均场态的对称性"时，我们真正指的是"相应自旋波函数 $\Psi_{\text{spin}}(\{\alpha_i\})$ 的对称性"。相应的自旋波函数由平均场态投影到每格点一个费米子的子空间得到（见式 (9.1.12)）。

由于平均场拟设 $\chi_{ij}$ 不依赖自旋取向，平均场基态中每一个负能级都被一个自旋向上和一个自旋向下的费米子占据。因此，平均场态 $|\Psi^{(\chi_{ij})}_{\text{mean}}\rangle$ 在自旋旋转下不变；结果，投影得到的物理自旋波函数也在自旋旋转下不变。

然而，平均场拟设在沿 $x$ 方向平移一个晶格间距下并非不变。因此，平均场态 $|\Psi^{(\chi_{ij})}_{\text{mean}}\rangle$ 破坏平移对称性。这似乎暗示物理自旋波函数也破坏平移对称性。事实上，物理自旋波函数并不破坏平移对称性。原因在于，平移后的拟设 $(\chi_{i,i+\pmb{x}},\chi_{i,i+\pmb{y}})=(-\mathrm{i}\chi_1(-)^{i_x},\mathrm{i}\chi_1)$ 通过取 $\mathrm{e}^{\mathrm{i}\theta_i}=(-)^{i_x}$ 的规范变换 (9.1.11) 与原拟设 $(\chi_{i,i+\pmb{x}},\chi_{i,i+\pmb{y}})=(\mathrm{i}\chi_1(-)^{i_x},\mathrm{i}\chi_1)$ 相联系。因此，这两个拟设给出同一个物理自旋波函数，投影后的自旋波函数在平移下不变。$\pi$ 通量相是一个平移对称的自旋液体。

第二，我们考虑低能激发的性质。半填充（即平均每格点一个费米子）时的费米面是位于 $(k_x,k_y)=(0,0)$ 和 $(0,\pi)$ 的点。在平均场层次上，$\pi$ 通量态含有无能隙自旋激发，它们对应于跨费米点的粒子–空穴激发。除这些无能隙自旋激发之外，平均场通量态还含有 $U(1)$ 规范涨落 $(a_{ij},a_0)$。在一阶平均场理论中，低能激发由耦合到 $U(1)$ 规范场的无能隙自旋子描写。

如果我们忽略自旋子 $f_i$ 与规范涨落 $(a_{ij},a_0)$ 之间的相互作用，那么 $\pi$ 通量相将具有由费米型准粒子描写的无能隙中性自旋 $1/2$ 激发。然而，除非证明了在计入规范相互作用后这一结果仍然成立，我们不应相信这个惊人的结果。

现在让我们考虑自旋激发与规范涨落之间的相互作用。为方便起见，让我们在哈密顿量 (9.1.10) 中加上 Maxwell 项 $\frac{1}{8\pi g^{2}}(v^{2}f_{12}^{2}-f_{0i}^{2})$，其中 $g$ 是耦合常数，$f_{\mu\nu}$ 是 $a_{\mu}$ 的场强。原理论对应于 $g\to\infty$ 的极限。Maxwell 项是在积掉费米子的过程中产生的。设 $g(\Lambda)$ 是积掉能标 $\Lambda$ 与 $\Lambda_0$ 之间的费米子后得到的耦合常数，其中 $\Lambda_0$ 是原理论中的能量截断。由量纲分析可得 $\mathrm{d}g^{-2}(\Lambda)\sim\mathrm{d}\Lambda/\Lambda^{2}$（注意 $\mathrm{d}g^{-2}(\Lambda)\propto\mathrm{d}\Lambda$）。于是 $g^{-2}(\Lambda)\sim\Lambda^{-1}-\Lambda_0^{-1}$。

由于与规范场耦合，一个自旋子会产生规范场 $(a_0,a_{ij})$ 的"电"场 $f_{0i}$（注意自旋子携带规范场的一个单位电荷）。粒子–空穴激发中粒子–空穴对之间的势能为
\begin{equation*}
V(\pmb{r}_1-\pmb{r}_2)=2\pi g^{2}\ln|\pmb{r}_1-\pmb{r}_2|
\tag{9.1.20}
\end{equation*}

我们发现粒子与空穴在长距离上仍有相互作用。为估计这一相互作用的强度和效应，让我们比较粒子–空穴对的动能 $E_K\sim 1/|\pmb{r}_1-\pmb{r}_2|$ 与势能 $V(\pmb{r}_1-\pmb{r}_2)$。由于 $g^{2}\sim v/|\pmb{r}_1-\pmb{r}_2|$（假定 $\Lambda\sim v/|\pmb{r}_1-\pmb{r}_2|$），我们看到 $V/E_K\sim 1$。这意味着相互作用是边缘的，在低能下不能忽略。这一结果提示，$\pi$ 通量态中的量子涨落（例如规范涨落）是重要的，将剧烈改变零阶平均场理论的低能性质。在这种情况下，零阶平均场理论（具有自由自旋子激发）不能给出自旋液体低能性质的可靠图像。要得到 $\pi$ 通量态可靠的低能性质，必须处理无能隙费米子与规范场的耦合系统。忽略规范相互作用而得到的那些惊人性质可能并不成立。\footnote{当平均场涨落微弱时，边缘平均场态可以导致代数自旋液体（algebraic spin liquid）——一种在低能下没有自由准粒子的自旋液体。见 9.9.5 节和 9.10 节。}

我们看到，投影构造的怀疑者们是对的：零阶平均场理论的结果有误导性，不可轻信。然而，相信者们也是对的：一阶平均场理论确实告诉我们，零阶平均场理论的结果是不正确的。到目前为止，投影构造没有误导过我们；只要计入适当的涨落，它是可以信赖的。与自旋液体物理性质相对应的，是一阶平均场理论（而非零阶平均场理论）的结果。

为了刻画平均场态周围涨落的重要性，这里我们想引入稳定、边缘和不稳定平均场态三个概念。在稳定平均场态中，涨落是微弱的，涨落诱导的相互作用在低能下消失（即相互作用是不相关微扰）。在边缘平均场态中，相互作用与能量之比在低能极限下趋于有限常数，此时相互作用是边缘微扰。在不稳定平均场态中，相互作用与能量之比在低能极限下发散，此时相互作用是相关微扰。只有对稳定平均场态，零阶平均场理论的性质才能在涨落之下幸存。对于不稳定平均场态，涨落会在低能下驱动相变，我们无法从零阶平均场理论推出自旋液体的任何物理性质。在这种情况下，一阶平均场理论也没有用，因为它并不能帮助我们推出自旋系统的物理性质。上面讨论的 $\pi$ 通量态是一个边缘平均场态。对于边缘平均场态，如果相互作用与能量之比很小，一阶平均场理论可能有用；此时相应自旋态的物理性质可以微扰地计算。对 $\pi$ 通量态而言，相互作用与能量之比是 1 的量级，因此很难从一阶平均场理论得到 $\pi$ 通量态的低能物理性质。

由上述讨论可见，一阶平均场理论只对稳定平均场态以及涨落微弱的边缘平均场态有用。用平均场理论研究自旋液体的关键，在于找到稳定平均场态（见 9.1.6、9.2.4 和 9.2.6 节），或涨落微弱的边缘平均场态（见 9.8 节）。

\subsection*{问题 9.1.2}

$\pi$ 通量态的旋转对称性

\begin{enumerate}
\item 证明式 (9.1.14)。
\item 证明 $\pi$ 通量态不破坏 $90^{\circ}$ 旋转对称性。
\end{enumerate}
